\documentclass[12pt]{article}
\usepackage[T1]{fontenc}
\usepackage[utf8]{inputenc}
\usepackage{lmodern}
\usepackage{textcomp}
\usepackage{enumitem}
\usepackage{geometry}
\geometry{margin=1in}
\begin{document}
\begin{center}
Answer Key - Chemistry - Graduate Level Assessment 13 \
\small (Answers are ordered exactly as in the question paper.)
\end{center}

\section*{Multiple Choice}
\begin{enumerate}[label=\arabic*.]
\item \textbf{Answer:} E
\\ \textit{Options:} A) 4.5 g; B) 2.2 g; C) 4.0 g; D) 2.0 g; E) 2.5 g
\\ \textit{Note:} The correct answer of 2.5 g of KCLO\_3 is derived from the stoichiometric decomposition reaction of potassium chlorate (KCLO\_3) into potassium chloride (KCl) and oxygen (O\_2), where the molar mass of KCLO\_3 and the molar ratio of oxygen produced indicate that 0.96 g of oxygen corresponds to the decomposition of 2.5 g of KCLO\_3 based on the balanced equation.
\item \textbf{Answer:} A
\\ \textit{Options:} A) CaS; B) MgS; C) NaCl; D) LiCl; E) NaF
\\ \textit{Note:} CaS has the greatest lattice enthalpy due to the combination of its high ionic charges ($Ca^2$+ and $S^2$-) and relatively small ionic radii, which lead to stronger electrostatic attractions between the ions.
\item \textbf{Answer:} A
\\ \textit{Options:} A) 2 mg \%; B) 0.0020 mg \%; C) 0.002 mg \%; D) 0.02 mg \%; E) 20 mg \%
\\ \textit{Note:} The correct answer is 2 mg \% because 1 ppm is equivalent to 1 mg of solute per liter of solution, and since 2.0 ppm translates directly to 2 mg of fluoride ion per 100 mL of water, it is expressed as 2 mg \%.
\item \textbf{Answer:} A
\\ \textit{Options:} A) 13.00573 amu; B) 13.00128 amu; C) 13.00987 amu; D) 13.00335 amu; E) 13.00739 amu
\\ \textit{Note:} The nuclidic mass of $N^13$ is 13.00573 amu because it accounts for the mass of the nucleus, including the mass defect due to the energy released in the beta-plus decay, which corresponds to the kinetic energy of the emitted positron and neutrino.
\item \textbf{Answer:} A
\\ \textit{Options:} A) 135 K; B) 150 K; C) 75 K; D) 110 K; E) 180 K
\\ \textit{Note:} The final temperature of the gas is 135 K because, during an adiabatic and reversible expansion, the temperature of an ideal gas decreases as the pressure decreases, which can be calculated using the relation T$_{2}$/T$_{1}$ = (P$_{2}$/P$_{1}$$)^($γ-1), where γ = C\_P/C\_V = 5/3 for a monatomic ideal gas.
\end{enumerate}

\section*{True / False}
\begin{enumerate}[label=\arabic*.]
\setcounter{enumi}{5}
\item \textbf{Answer:} False
\\ \textit{Options:} A) True; B) False
\\ \textit{Note:} The statement is false because a half-life that varies with concentration suggests a reaction order other than second order, as second-order reactions have a half-life that is inversely proportional to the initial concentration.
\item \textbf{Answer:} True
\\ \textit{Options:} A) True; B) False
\\ \textit{Note:} The hydronium ion concentration of a 0.1 M acetic acid solution is approximately 0.0013 M because acetic acid is a weak acid that partially dissociates in solution, leading to a much lower concentration of hydronium ions than the initial acid concentration.
\item \textbf{Answer:} False
\\ \textit{Options:} A) True; B) False
\\ \textit{Note:} The magnitude of the spin magnetic moment of an electron is actually 9.27 × 10^\{-24\} J/T, which is significantly lower than the stated value of 2.20 × 10^\{-23\} J/T.
\item \textbf{Answer:} True
\\ \textit{Options:} A) True; B) False
\\ \textit{Note:} Transition metals typically lose their outermost s-block electrons first due to their electron configuration, which facilitates the formation of bonds and results in a common oxidation state of +2.
\item \textbf{Answer:} False
\\ \textit{Options:} A) True; B) False
\\ \textit{Note:} The wave function for a 2 P\_Z electron varies with angle because it is dependent on the spherical harmonics, which describe the angular distribution of probability density in three-dimensional space.
\end{enumerate}

\section*{Long Form Response}
\begin{enumerate}[label=\arabic*.]
\setcounter{enumi}{10}
\item \textbf{Answer:} The principal allotropes of elemental boron, including alpha-boron (α-B) and beta-boron (β-B), exhibit distinct solid-state structures characterized by the presence of B$_{5}$ pentagons, which play a crucial role in their stability and properties. In α-B, these pentagons are arranged in a three-dimensional network, contributing to the material's hardness and high melting point, while also influencing its electronic properties. Conversely, β-B features a more complex arrangement of B$_{12}$ icosahedra and B$_{5}$ pentagons, enhancing its structural diversity and leading to different physical characteristics. The significance of B$_{5}$ pentagons lies in their ability to form stable, interconnected frameworks that dictate the overall geometry and bonding, ultimately affecting the mechanical and thermal properties of boron allotropes. Thus, the role of B$_{5}$ pentagons is pivotal in understanding the unique characteristics of boron's allotropes and their potential applications in advanced materials.
\\ \textit{Note:} correct because it accurately describes the distinct solid-state structures of the principal allotropes of boron, emphasizing the significance of B$_{5}$ pentagons in α-B for stability and properties, while also noting the structural differences in β-B that contribute to its unique characteristics.
\item \textbf{Answer:} To calculate the heat energy required to raise the temperature of 3.0 liters of water from 20^\ensuremath{\circC} to 80^\ensuremath{\circC}, we can use the formula \( Q = mc\Delta T \), where \( Q \) is the heat energy, \( m \) is the mass of the water, \( c \) is the specific heat capacity of water (approximately 1 kcal/kg·^\ensuremath{\circC}), and \( \Delta T \) is the change in temperature. Given that the density of water is about 1 kg/L, the mass of 3.0 liters of water is 3.0 kg. The temperature change \( \Delta T \) is 60^\ensuremath{\circC} (80^\ensuremath{\circC} - 20^\ensuremath{\circC}). Substituting these values into the equation gives \( Q = 3.0 \, \text{kg} \times 1 \, \text{kcal/kg·^\circC} \times 60^\circC = 180 \, \text{kcal} \). This calculation illustrates the principles of heat transfer and specific heat capacity, which describe how much energy is needed to change the temperature of a substance based on its mass and inherent properties.
\\ \textit{Note:} correct because it accurately applies the formula for heat energy, incorporating the mass of the water, its specific heat capacity, and the temperature change, which collectively demonstrate the principles of heat transfer and specific heat capacity in a clear and systematic manner.
\end{enumerate}

\vfill
\noindent\textit{End of Answer Key}
\end{document}